\chapter{L'application}
\label{chap:application}

\begin{chapeau}
Ce chapitre explique comment les étapes de calcul deviennent une application de bureau : les couches du logiciel, le
serveur local, l'interface, la fenêtre, l'application autonome, comment la lancer sur Mac et sur Windows, et les
fichiers qu'elle produit.
\end{chapeau}

\section{Des étapes de calcul à une application}

Les fichiers \fichier{etape1} à \fichier{etape7} restent des modules Python ordinaires, utilisables en ligne de
commande. L'application ajoute trois couches autour d'eux, sans les réécrire (\cref{fig:architecture}) :
\begin{enumerate}
  \item un \textbf{serveur local} (\fichier{application/serveur.py}) qui importe les modules et les appelle à la
  demande ;
  \item une \textbf{interface web} (HTML et JavaScript) qui envoie les clics de l'utilisateur au serveur et affiche les
  résultats ;
  \item une \textbf{fenêtre native} (\fichier{application/fenetre.py}) qui affiche cette interface comme une vraie
  application, avec son nom et son icône.
\end{enumerate}
Enfin, \textbf{PyInstaller} peut emballer le tout dans un seul programme, avec Python et toutes les bibliothèques.

\fig{architecture_application.pdf}{L'organisation du logiciel : tout reste sur l'ordinateur.}{fig:architecture}

\section{Le serveur local}
\label{sec:serveur}

\subsection{Le principe}

Le serveur est un petit serveur web, construit avec le module \code{http.server} de Python. Il n'écoute que l'adresse
\code{127.0.0.1} : il n'est joignable que depuis l'ordinateur lui-même, et rien ne passe par Internet. Il choisit un
port libre au démarrage. Chaque bouton de l'interface correspond à une adresse (\cref{tab:api}).

\begin{table}[htbp]
\centering\small
\caption{Les principales adresses du serveur.}
\label{tab:api}
\begin{tabularx}{\linewidth}{L{4.2cm}Y}
\toprule
\titrecol{Adresse} & \titrecol{Rôle} \\
\midrule
\code{/api/action} & lance un calcul : \code{inventaire}, \code{ouvrir}, \code{segmenter}, \code{nommer},
\code{separer}, \code{oublier}, \code{guide} \\
\code{/api/etat} & donne l'avancement du calcul en cours (étape, journal, erreur, résultat) \\
\code{/api/serie}, \code{/api/coupe} & décrit la série ouverte ; renvoie une coupe du scanner en image, os colorés \\
\code{/api/point}, \code{/api/maillage} & valeur sous le curseur ; maillage 3D d'un os \\
\code{/api/cora/cas}, \code{/api/cora/os}, \code{/api/cora/maillage} & l'os à planifier, ses repères et son maillage dans le repère de l'os \\
\code{/api/cora/guide}, \code{/api/cora/guide/fichier} & le résultat de l'étape 7 et ses fichiers (affichage, téléchargement) \\
\code{/api/choisir_dossier}, \code{/api/montrer} & choisir un dossier de scanners ; ouvrir le dossier des résultats \\
\code{/api/quitter} & arrêter le serveur \\
\bottomrule
\end{tabularx}
\end{table}

\subsection{Un calcul à la fois, dans le fil principal}

Sur macOS, la bibliothèque VTK exige de tourner dans le fil principal du programme. Les calculs y sont donc exécutés,
un à la fois ; le serveur HTTP tourne à côté et ne fait que lire (coupes, maillages, fichiers). Quand l'utilisateur
clique sur \ui{Segmenter la série}, la demande est mise en file, le serveur répond tout de suite, et l'interface
interroge \code{/api/etat} plusieurs fois par seconde pour afficher l'avancement (\cref{fig:parcours}). Pendant un
calcul, le serveur empêche la mise en veille de l'ordinateur par inactivité.

\fig[0.72]{parcours_requete.pdf}{Le parcours d'une demande, du clic au résultat.}{fig:parcours}

Par exemple, la fonction \code{segmenter()} du serveur applique les réglages choisis dans l'interface (seuil, lissage)
aux paramètres de \fichier{parametres.py}, appelle \code{lancer_tout.traiter()}, c'est-à-dire \textbf{les mêmes étapes
1 à 6 qu'en ligne de commande}, puis remet les paramètres d'origine et garde le résultat en mémoire.

\section{L'interface}

\subsection{Onglet 1 · Segmentation CT}

La liste des séries trouvées, les coupes du scanner (axiale, coronale, sagittale) avec les os en couleur, les os en 3D
et la liste des os avec leur nom, leur volume et la raison de ce nom (\cref{fig:capture-segmentation}). On peut
\ui{Renommer} un os ou le \ui{Séparer} en cliquant des points. La carte \ui{Os à planifier (CORA)} propose l'os complet
de la série, avec un bouton \ui{Planifier CORA $\rightarrow$}.

\fig{capture_segmentation.png}{L'onglet \ui{Segmentation CT} avec une série segmentée : coupes du scanner, os en 3D et
liste des os trouvés.}{fig:capture-segmentation}

\subsection{Onglet 2 · Planification CORA}

Deux sources : l'\textbf{os segmenté} (celui de la série ouverte) ou l'\textbf{os fictif}. Un panneau à gauche permet
de choisir la référence (valeurs normales ou os avant déformité), une déformité de test, un des 17 préréglages
cliniques, la stratégie (type de coin, niveau du trait, profondeur de la charnière, broches) et les 27 dimensions du
guide. Au centre, la vue 3D (Three.js) ; à droite, le texte de l'étape en cours, parmi les 13 étapes
(\cref{tab:13etapes}). À l'étape 12, le bouton \ui{Construire les guides 3D} envoie la prescription au serveur, qui
lance l'étape 7.

\begin{pourquoi}
\textbf{Pourquoi un seul fichier HTML ?} Les modules JavaScript de \fichier{code_source/cora/app/} sont assemblés par
\fichier{code_source/cora/build_app.py} en un seul fichier, \fichier{application/interface/cora/index.html}, avec
Three.js, les polices et les données de l'os fictif. Il s'ouvre sans Internet et sans installation. En contrepartie,
il faut relancer \fichier{build_app.py} après chaque modification du JavaScript.
\end{pourquoi}

\subsection{La 3D à la demande}

Les vues 3D ne sont redessinées que quand quelque chose change (rotation, curseur, étape). Au repos, l'application ne
calcule aucune image : le processeur reste libre et l'ordinateur ne chauffe pas.

\section{La fenêtre native}

\fichier{application/fenetre.py} utilise la bibliothèque \textbf{pywebview}, qui affiche une page web dans une fenêtre
du système : le moteur WebKit (celui de Safari) sur macOS, WebView2 (celui d'Edge) sur Windows. L'application a ainsi
son nom et son icône dans le Dock ou la barre des tâches, des menus en français et des boîtes « Enregistrer sous » pour
les exports. Si pywebview ne peut pas être installé, l'application s'ouvre dans une fenêtre de Chrome ou d'Edge.

\section{L'application autonome}

Pour qu'un utilisateur n'ait rien à installer, PyInstaller copie Python, toutes les bibliothèques et le code dans une
application d'environ 800~Mo. Au lancement, un seul exécutable démarre \textbf{deux processus} : la fenêtre, et le
serveur de calcul (le même exécutable, lancé avec l'option \code{--serveur}). Sur macOS, la fenêtre (Cocoa) et les
calculs (VTK) exigent tous deux le fil principal : deux processus les séparent. Fermer la fenêtre arrête le serveur,
après le calcul en cours.

\begin{attention}
L'application autonome contient une \textbf{copie figée} du code : il faut la reconstruire après chaque modification
(\fichier{application/lanceurs/creer_app_autonome.sh} sur Mac, environ 3 minutes ;
\fichier{creer_app_windows.bat} sur Windows, environ 10 minutes la première fois).
\end{attention}

\section{Lancer le logiciel}
\label{sec:lancer}

\begin{table}[htbp]
\centering\small
\caption{Les façons de lancer CT-CORA3D.}
\label{tab:lancer}
\begin{tabularx}{\linewidth}{L{3.2cm}Y}
\toprule
\titrecol{Système} & \titrecol{Comment} \\
\midrule
macOS (conseillé) & \fichier{~/Applications/CT-CORA3D.app} : application autonome \\
macOS (développement) & \fichier{CT-CORA3D.app} du dossier du projet : petit lanceur qui utilise le Python de l'ordinateur \\
Windows (développement) & \fichier{Lancer sur Windows.bat} : essayé avec succès sur un PC Windows le 04/11/2026 \\
Windows (autonome) & raccourci créé par \fichier{application/lanceurs/creer_app_windows.bat} (recette écrite, pas encore construite) \\
\bottomrule
\end{tabularx}
\end{table}

Au premier lancement, les lanceurs de développement créent un environnement Python propre à l'ordinateur, hors de
OneDrive (\fichier{~/Library/Application Support/CTCORA3D/venv} sur Mac, \texttt{\%LOCALAPPDATA\%\textbackslash{}CTCORA3D\textbackslash{}venv} sur
Windows), y installent les bibliothèques (quelques minutes, Internet nécessaire), puis pywebview. Les lancements
suivants prennent quelques secondes : le lanceur vérifie seulement que les bibliothèques sont présentes, sans les
charger (0,02~s).

\subsection{Depuis GitHub}

Le code est versionné dans un dépôt GitHub \textbf{public} : chacun peut le récupérer
(\code{git clone}), installer Python 3.12 et lancer l'application comme ci-dessus. Les scanners DICOM ne sont pas dans le
dépôt (ils sont trop lourds) : il faut les placer dans le dossier \fichier{dicom/} de \fichier{CT-CORA3D/} (ou dans un
dossier \fichier{CT-scan/} à côté), ou les choisir par le bouton \ui{Changer de dossier…}. Un fichier \fichier{.gitattributes} impose les fins de ligne Windows aux fichiers
\fichier{.bat} : sans cela, un \fichier{.bat} récupéré depuis GitHub peut échouer sous Windows.

\section{Les fichiers produits}

\begin{table}[htbp]
\centering\small
\caption{Le contenu de \fichier{resultats/<série>/}.}
\label{tab:fichiers}
\begin{tabularx}{\linewidth}{L{4.6cm}Y}
\toprule
\titrecol{Fichier} & \titrecol{Contenu} \\
\midrule
\fichier{femur.stl}, \fichier{tibia.stl}\ldots & un maillage étanche par os, en millimètres, repère du scanner \\
\fichier{segmentation.seg.nrrd} & toutes les étiquettes, avec noms et couleurs (s'ouvre dans 3D Slicer) \\
\fichier{rapport.md}, \fichier{rapport.json} & le rapport de la segmentation \\
\fichier{apercu_3d.png}, \fichier{apercu_coupes.png} & des images de contrôle \\
\fichier{reperes_auto.mrk.json} & la tête du fémur, au format de 3D Slicer \\
\fichier{cora_<os>.json}, \fichier{.bin}, \fichier{.png} & repère, côté, repères, angles, torsion ; maillage dans le repère de l'os \\
\fichier{cora_cas.json} & l'os à planifier et la raison du choix \\
\fichier{guide_<os>_<coin>/} & \fichier{guide_coupe.stl} et \fichier{guide_reduction.stl} (à imprimer), \fichier{guide_coupe_scanner.stl}, \fichier{os_corrige.stl}, \fichier{prescription.json}, \fichier{guide.json} (contrôles, mode opératoire), plans 2D \\
\bottomrule
\end{tabularx}
\end{table}

Il n'y a pas de base de données : pour sauvegarder un cas, il suffit de copier son dossier. Le journal de
l'application est dans \fichier{~/Library/Logs/CTCORA3D.log} (Mac) ou \texttt{\%LOCALAPPDATA\%\textbackslash{}CTCORA3D\textbackslash{}CTCORA3D.log}
(Windows).

\section{Les performances}

\begin{table}[htbp]
\centering\small
\caption{Mesures sur le Mac de développement, avant et après le travail sur la fluidité (22/10/2026). Les résultats
des guides sont identiques au bit près.}
\label{tab:perf}
\begin{tabularx}{\linewidth}{Yrr}
\toprule
\titrecol{Mesure} & \titrecol{Avant} & \titrecol{Après} \\
\midrule
guide complet (tibia, fermeture, guide v3 avec verrou, plans 2D) & 86,5~s & 23,5~s \\
\quad dont plans 2D & 49,8~s & 6,2~s \\
chargement des bibliothèques du serveur & 2,1~s & 0,9~s \\
vérification de l'environnement par les lanceurs & 1,5 à 2,6~s & 0,02~s \\
vue 3D au repos (images calculées par seconde) & 60 & 0 \\
\bottomrule
\end{tabularx}
\end{table}

Les gains viennent de choix simples : plans 2D dessinés par un tampon de profondeur en numpy, recherches de voisins
sur tous les cœurs du processeur, et chargement des seules classes de VTK utilisées (\fichier{vtk_leger.py}). Depuis
la simplification du guide (06/11/2026), un guide complet avec ses plans prend 10 à 16~s.

\begin{danslecode}
\fichier{application/serveur.py} (serveur, actions, adresses), \fichier{application/fenetre.py} (fenêtre native),
\fichier{application/ct_cora3d_app.py} (point d'entrée de l'application autonome), \fichier{application/interface/}
(onglets), \fichier{application/lanceurs/} (fabrication des applications), \fichier{application/lancer_mac.sh} et
\fichier{Lancer sur Windows.bat} (lanceurs de développement).
\end{danslecode}
